\documentclass[12pt,a4paper]{article}

% ── Packages ──────────────────────────────────────────────
\usepackage[top=1.8cm, bottom=2cm, left=2cm, right=2cm]{geometry}
\usepackage{times}
\usepackage[utf8]{inputenc}
\usepackage{multicol}
\usepackage{graphicx}
\usepackage{amsmath}
\usepackage{amssymb}
\usepackage{xcolor}
\usepackage{enumitem}
\usepackage{tabularx}
\usepackage{fancyhdr}
\usepackage{lastpage}

% ── No paragraph indent, small paragraph skip ────────────
\setlength{\parindent}{0pt}
\setlength{\parskip}{4pt}

% ── Disable default page numbering ───────────────────────
\pagestyle{empty}

% ── Custom list styles (shared with literature template) ──
\newlist{romanenum}{enumerate}{1}
\setlist[romanenum]{label=\textbf{(\roman*)},leftmargin=1.2cm,itemsep=3pt,topsep=4pt}

\newlist{abcenum}{enumerate}{1}
\setlist[abcenum]{label=\textbf{(\alph*)},leftmargin=1.5cm,itemsep=1pt,topsep=2pt}

% ── Marks flush-right ────────────────────────────────────
\newcommand{\qmarks}[1]{\hfill\textbf{[#1]}}

% ══════════════════════════════════════════════════════════
%          CUSTOM MACROS (API unchanged from v1)
% ══════════════════════════════════════════════════════════

% ── Header ────────────────────────────────────────────────
% \hdr{School Name}{Address}{Subject}{Class -- Exam Type}{Time: N min}{Maximum Marks: N}
\newcommand{\hdr}[6]{%
\begin{center}
    {\LARGE\bfseries #1}\\[2pt]
    {#2}\\[6pt]
    {\Large\bfseries #3}\\[2pt]
    {\Large\bfseries #4}
\end{center}
\vspace{4pt}
\noindent\rule{\textwidth}{0.8pt}
\vspace{-6pt}
\noindent\begin{tabularx}{\textwidth}{@{}X r@{}}
\textbf{\textit{#5}} & \textbf{\textit{#6}}
\end{tabularx}
\vspace{-2pt}
\noindent\rule{\textwidth}{0.4pt}
\vspace{4pt}}

% ── General Instructions ─────────────────────────────────
% \geninstructions{\item ...\item ...}
\newcommand{\geninstructions}[1]{%
\textbf{General Instructions:}
\begin{enumerate}[leftmargin=0.8cm, itemsep=1pt, topsep=2pt]
    #1
\end{enumerate}
\vspace{2pt}
\noindent\rule{\textwidth}{0.4pt}
\vspace{6pt}}

% ── Question Section ─────────────────────────────────────
% \questionsection{A}{MCQs (1 mark each)}{section instructions as \item list}
\newcommand{\questionsection}[3]{%
\vspace{8pt}
\begin{center}
    {\large\bfseries SECTION~--~#1}\\[2pt]
    \textbf{#2}
\end{center}
\vspace{2pt}
{\small\textit{%
\begin{itemize}[noitemsep, topsep=0pt, leftmargin=0.8cm]
    #3
\end{itemize}}}
\vspace{4pt}}

% ── 4-Option MCQ — 2-column layout ───────────────────────
% \mcqfourtwo{num}{text}{optA}{optB}{optC}{optD}{(N mark)}{hint}{extra}
\newcommand{\mcqfourtwo}[9]{%
\vspace{4pt}
\textbf{Q#1.}~#2 \hfill \textbf{#7}\\[2pt]
\ifx&#8&\else{\small\textit{#8}}\\[1pt]\fi
#9
\vspace{2pt}
\begin{tabularx}{\textwidth}{@{\hspace{1.2cm}}X X}
\textbf{(a)}~#3 & \textbf{(b)}~#4 \\
\textbf{(c)}~#5 & \textbf{(d)}~#6
\end{tabularx}
\vspace{2pt}}

% ── 4-Option MCQ — 1-column layout (for long options) ────
% \mcqfour{num}{text}{optA}{optB}{optC}{optD}{(N mark)}{hint}{extra}
\newcommand{\mcqfour}[9]{%
\vspace{4pt}
\textbf{Q#1.}~#2 \hfill \textbf{#7}\\[2pt]
\ifx&#8&\else{\small\textit{#8}}\\[1pt]\fi
#9
\begin{abcenum}
    \item #3
    \item #4
    \item #5
    \item #6
\end{abcenum}
\vspace{2pt}}

% ── Descriptive / Fill / True-False Question ─────────────
% \question{num}{text}{(N marks)}{hint}
\newcommand{\question}[4]{%
\vspace{3pt}
\textbf{Q#1.}~#2 \hfill \textbf{#3}%
\ifx&#4&\else\\[1pt]{\small\textit{#4}}\fi
\\[2pt]}

% ── Centered equation helper ─────────────────────────────
\newcommand{\equ}[1]{%
\begin{center}$\displaystyle #1$\end{center}}

% ══════════════════════════════════════════════════════════
%              BEGIN DOCUMENT CONTENT BELOW
% ══════════════════════════════════════════════════════════

\begin{document}

% ── Example header ────────────────────────────────────────
\hdr{Delhi Public School}{R.K.\ Puram, New Delhi}{Science -- Forces and Motion}{Class VIII -- Unit Test}{Time: 45 minutes}{Maximum Marks: 30}

% ── General instructions ──────────────────────────────────
\geninstructions{
\item All questions are compulsory.
\item Read each question carefully before answering.
\item Write neat and legible answers.
\item Marks are indicated against each question in brackets~[~].
}

% ══════════════════════════════════════════════════════════
%                       SECTION A
% ══════════════════════════════════════════════════════════

\questionsection{A}{Multiple Choice Questions (1 mark each)}{
\item Choose the correct option for each question.
\item Each question carries 1 mark. No negative marking.}

\mcqfourtwo{1}{A force can change the:}
    {colour of an object}
    {shape of an object}
    {smell of an object}
    {material of an object}
    {[1]}{}{}

\mcqfourtwo{2}{The SI unit of force is:}
    {Joule}
    {Watt}
    {Newton}
    {Pascal}
    {[1]}{}{}

\mcqfourtwo{3}{Friction can be reduced by:}
    {increasing the weight}
    {making surfaces rough}
    {using lubricants}
    {increasing contact area}
    {[1]}{}{}

\mcqfourtwo{4}{Gravity is a:}
    {contact force}
    {non-contact force}
    {frictional force}
    {muscular force}
    {[1]}{}{}

\mcqfourtwo{5}{Which of the following is a contact force?}
    {Magnetic force}
    {Gravitational force}
    {Electrostatic force}
    {Frictional force}
    {[1]}{}{}

% ══════════════════════════════════════════════════════════
%                       SECTION B
% ══════════════════════════════════════════════════════════

\questionsection{B}{Fill in the Blanks (1 mark each)}{
\item Fill in the blanks with the correct word or phrase.}

\question{6}{The force exerted by a magnet is called \underline{\hspace{3cm}} force.}{[1]}{}
\question{7}{The force of \underline{\hspace{3cm}} acts between all objects in the universe.}{[1]}{}
\question{8}{A force can change the \underline{\hspace{3cm}} of motion of a body.}{[1]}{}

% ══════════════════════════════════════════════════════════
%                       SECTION C
% ══════════════════════════════════════════════════════════

\questionsection{C}{Short Answer Questions (2 marks each)}{
\item Answer in 2--3 sentences.}

\question{9}{Define force and state its SI unit.}{[2]}{}
\question{10}{Differentiate between contact and non-contact forces with one example each.}{[2]}{}
\question{11}{Why do we sprinkle fine powder on a carrom board?}{[2]}{}

% ══════════════════════════════════════════════════════════
%                       SECTION D
% ══════════════════════════════════════════════════════════

\questionsection{D}{Long Answer Questions (4 marks each)}{
\item Answer in 5--6 sentences with appropriate detail.}

\question{12}{Explain the different types of friction. Which type of friction is the least?}{[4]}{}
\question{13}{Describe the effects of force on an object. Give two examples for each effect.}{[4]}{}

% ══════════════════════════════════════════════════════════
%                      ANSWER KEY
% ══════════════════════════════════════════════════════════

\newpage

\begin{center}
\textbf{\Large ANSWER KEY}\\[3pt]
\rule{0.5\textwidth}{0.4pt}
\end{center}

\vspace{6pt}

\textbf{\underline{Section A -- Multiple Choice Questions}}\\[3mm]
\begin{tabular}{l l}
Q1. & (b) shape of an object \\
Q2. & (c) Newton \\
Q3. & (c) using lubricants \\
Q4. & (b) non-contact force \\
Q5. & (d) Frictional force \\
\end{tabular}

\vspace{6pt}
\textbf{\underline{Section B -- Fill in the Blanks}}\\[3mm]
\begin{tabular}{l l}
Q6. & magnetic \\
Q7. & gravity (gravitational) \\
Q8. & speed / direction \\
\end{tabular}

\vspace{6pt}
\textbf{\underline{Section C -- Short Answers (Key Points)}}\\[3mm]
\textbf{Q9.} A push or pull acting on an object is called a force. SI unit: Newton (N).\\[2mm]
\textbf{Q10.} Contact forces require physical contact (e.g.\ friction). Non-contact forces act at a distance (e.g.\ gravity).\\[2mm]
\textbf{Q11.} Fine powder reduces friction, allowing the striker to slide smoothly.\\[2mm]

\vspace{6pt}
\textbf{\underline{Section D -- Long Answers (Key Points)}}\\[3mm]
\textbf{Q12.} Static friction (body at rest), sliding friction (body in motion), rolling friction (rolling objects). Rolling friction is the least.\\[2mm]
\textbf{Q13.} Force can change shape (squeezing clay), speed (pushing a cart), direction (hitting a ball), or start/stop motion (kicking a ball, applying brakes).\\[2mm]

\vspace{20pt}
\begin{center}
$\blacksquare$~$\blacksquare$
\end{center}

\end{document}
